\chapter{Generators and mesh I/O}
\label{ch:generators}

\section{Generators}

All generators fill an \emph{empty} object that has been added to an ensemble. Rank 0 builds
the whole mesh, the id counters are broadcast, and the other ranks stay empty until the caller
redistributes. The coordinates given to a generator are \emph{continuous}: a centre outside the
box, or a body that crosses a periodic face, is wrapped into $(\texttt{pos},\texttt{img})$ and
canonicalised. The closed generators wind outward ($V>0$).

\begin{lstlisting}[numbers=none]
int  MFTLGenerateSphere(MFTL_Object *object, double center[3], double radius, int refinements);
void MFTLGenerateCylinder(MFTL_Object *obj, double dr, double radius, double L, double ellipticity);
void MFTLGenerateBox(MFTL_Object *obj, const double lo[3], const double hi[3]);
void MFTLGenerateCube(MFTL_Object *obj, int refine);
void MFTLGenerateFace(MFTL_Object *obj, const double v[4][3]);
\end{lstlisting}
\begin{description}
  \item[Sphere] an icosahedron refined \texttt{refinements} times by 1-to-4 subdivision with the
    new vertices projected onto the sphere: $20\cdot4^{r}$ triangles (1280 for $r=3$). The call is
    collective (it checks the minimum-image constraint).
  \item[Cylinder] a capped cylinder along $+z$, $z\in[0,L]$, with the \mftlone{}
    parameterisation of the resolution by \texttt{dr}; \texttt{ellipticity} in $[0,1)$ scales the
    $x$ semi-axis by $\sqrt{1-e^2}$, a value in $(-1,0)$ the $y$ semi-axis.
  \item[Box, cube] closed axis-aligned boxes (the cube is $[-1,1]^3$ refined by flat 1-to-4
    subdivision).
  \item[Face] an open quadrilateral patch of two triangles, declared open with one boundary loop
    (tag 0).
\end{description}
After generating, call \fn{MFTLComputeTriangleProperties} and redistribute
(\S\ref{sec:ex-basic}).

\section{GTS surface files}

\begin{lstlisting}[numbers=none]
int MFTLReadGTS(MFTL_Object *obj, const char *path);    /* collective */
int MFTLWriteGTS(MFTL_Object *obj, const char *path);
\end{lstlisting}
\fn{MFTLReadGTS} parses the \mftlone{} GTS layout (a header with the numbers of vertices, edges
and triangles, the vertex coordinates, 1-based edge vertex pairs and triangles as 1-based edge
triples) on rank 0 and builds the object there. The triangles are oriented consistently, and a
closed body is canonicalised to outward winding. The coordinates are continuous input. An edge
with a single triangle is accepted only on a body declared open before the call; its loops
become tagged boundaries. \fn{MFTLWriteGTS} writes the triangles owned by rank 0 with
\texttt{\%.17g} coordinates (a bit-exact round trip); it refuses with \texttt{MFTL\_ERROR} when the
object has been distributed (call it before redistributing, or inside a serial bracket).

\section{Visualisation output}

\fn{MFTLWriteVtk}\texttt{(obj, base)} writes the entities stored on the calling rank (owned and
ghost) to a legacy VTK file \texttt{base.vtk}. \fn{MFTLWriteVtu}\texttt{(ens, prefix)} and its
alias \fn{MFTLWritePvtu} write all objects as a parallel VTK set: one \texttt{prefix\_NNNNN.vtu}
piece per rank and a \texttt{prefix.pvtu} index on rank 0, with the vertex owner, the triangle
owner and the body id as data. A triangle that crosses a periodic seam is written with its own
copies of its corners in the image of its first corner, so it is not stretched across the box.

\section{Native restart meshes}

\begin{lstlisting}[numbers=none]
int MFTLWriteNativeMesh(const MFTL_Object *obj, const char *path);    /* v1 (MFTL1 .mesh) */
int MFTLWriteNativeMeshV2(const MFTL_Object *obj, const char *path);  /* v2 */
int MFTLReadNativeMesh(MFTL_Object *obj, const char *path);           /* v1 or v2, collective */
\end{lstlisting}
The writers write the entities stored on the calling rank, so an object is first gathered on
one rank (\fn{MFTLEnsembleBeginSerial}, \S\ref{sec:serial-bracket}). Version~1 is the
\mftlone{} \texttt{.mesh} format with the surface scalar. Version~2 adds the declared class, the
period box, exact (\texttt{\%.17g}) wrapped positions with their image counters, the boundary
tags and explicit vertex triples. \fn{MFTLReadNativeMesh} builds an empty object from either
version on any rank count (ranks that pass \texttt{NULL} contribute nothing); call
\fn{MFTLRedistribute} afterwards. The reader re-establishes the canonical image gauge, so a
state written in the canonical gauge is restored exactly --- \multi{} restarts from version~2
files.
